\section{Introduction}
\label{sec:intro}

Spatial transcriptomics (ST) measures gene expression at spatially
resolved positions in tissue \citep{stahl2016visualization,
chen2015spatially,wang2018three}. The dominant sequencing-based
platforms (e.g., 10x Genomics Visium, Slide-seq) measure expression
at \emph{spots} that aggregate $\sim 1$--$50$ cells of unknown type
composition; the imaging-based platforms (MERFISH, seqFISH+) achieve
single-cell or subcellular resolution but for a limited gene set.
Inferring per-cell-type expression from spot-level data is the
\emph{deconvolution problem}, which has spawned a rapidly growing
methods literature: Cell2location \citep{kleshchevnikov2022cell2location},
RCTD \citep{cable2021robust}, Tangram \citep{biancalani2021deep},
CARD \citep{ma2022spatially}, SPOTlight \citep{elosua2021spotlight},
Stereoscope \citep{andersson2020stereoscope},
SpatialDecon \citep{danaher2022advances},
DestVI \citep{lopez2022destvi},
and the reference-free STdeconvolve \citep{miller2022reference},
among the published methods benchmarked by \citet{li2023comprehensive}.

These methods solve the deconvolution problem with a variety of
parametric and non-parametric approaches. Several discuss
identifiability informally (RCTD via its likelihood structure;
Cell2location via its prior; CARD via spatial-smoothness
restoration), but no prior work characterizes deconvolution
identifiability through the lens of masked-data coarsening
conditions, and the consequences for method choice across the
literature have not been mapped. Practical questions go unanswered:
how many marker genes are needed, and when can deconvolution fail
silently?

\subsection{The bridge to masked-data inference}

We observe that ST deconvolution is mathematically isomorphic to the
masked-data series-system identifiability problem of
\citet{towell2026masked}. In that framework, $m$ components in
series produce a system that fails when any one fails; the failed
component is masked but a \emph{candidate set} containing it is
observed. Identifiability conditions on the candidate-set matrix
determine when component-level parameters can be recovered.

The translation to ST is direct: cell types play the role of
components, spots play the role of candidate sets (observed cell-type
composition vectors), single-cell-resolution platforms provide
singleton candidate sets, and marker-gene panels play the role of
identifying functions. The C1--C2--C3 conditions for ignorable
coarsening port directly. This framework was previously applied to
the related scRNA-seq zero-inflation problem
\citep{towell2026scrnacoarsening,jiang2022zero}; ST is a strict
generalization in that the candidate set has non-trivial structure
(multiple cell types per spot) rather than being a single bit (zero
vs. nonzero count).

\subsection{Contributions}

We contribute:

\begin{itemize}
\item \textbf{A bridge from ST deconvolution to masked-data
  inference} (\cref{sec:translation}), formalizing existing methods
  (RCTD, Cell2location, Tangram, CARD) as special cases of the
  framework.

\item \textbf{An identifiability theorem} (\cref{sec:identifiability}):
  cell-type-specific expression is identifiable from spot-level data
  if and only if the spot-by-cell-type composition matrix has full
  column rank. We discuss when this condition is met empirically, and
  when single-cell-resolution platforms or marker-gene panels are
  needed to restore it.

\item \textbf{A cell-total consistency theorem}
  (\cref{sec:identifiability}): naive deconvolution methods exactly
  match observed spot-level expression in aggregate, regardless of
  per-cell-type bias. This recovers the analogue of the scRNA-seq
  result of \citet{towell2026scrnacoarsening}.

\item \textbf{A bias bound under marker-gene undercoverage}
  (\cref{sec:methodology}): a closed-form first-order Taylor bound on
  the deconvolution bias as a function of the gap between the chosen
  marker-gene set and the ``true'' identifying set. Recovers as a
  corollary Tangram's observation that ``a few hundred marker genes
  suffice'' in mouse brain cortex \citep{biancalani2021deep}.

\item \textbf{Simulation and real-data validation}
  (\cref{sec:validation}). On simulated data, cell-total consistency
  holds to machine precision in the noiseless limit and to the
  Poisson noise floor at finite spot count, and the rank condition is
  verified across $200$ replicates. On the public 10x Visium adult
  mouse coronal brain dataset, NMF at $K = 20$ explains $92\%$ of
  the variance in the log-normalized expression matrix, and a
  restricted-region scenario reproduces the singleton-augmentation
  mechanism the framework predicts.
\end{itemize}

The paper's central message: ST deconvolution methods are not
arbitrary heuristics; they are special cases of a single
identifiability framework that predicts when each method is
appropriate. Practitioners can use the rank condition as a diagnostic
before applying any deconvolution method.
